\chapter{Open Questions and Research Directions}
\label{ch:open-questions}

\epigraph{A working document that hides its weaknesses is less useful than one that maps them. These gaps are research directions, not admissions of failure.}{}

This chapter is an exercise in intellectual honesty. The coordination collapse thesis, as presented in this report, is incomplete. The evidence base has specific weaknesses. The analysis contains blind spots visible to each of its three audiences. Several important topics are missing entirely. Mapping these gaps is not an admission of failure. It is a research agenda.

\section{The Core Thesis Is Incomplete}

The report frames hierarchy as an information routing protocol and argues that AI collapses the cost of that routing to near-zero. This framing is analytically productive: it generates the coordination collapse concept, the judgment broker role, and the compounding organisation model. But hierarchy performs at least five functions, of which information routing is only one.

\subsection{Legitimation}

Who has authority to commit resources? The mesh can route information and even recommend resource allocation, but it cannot \emph{authorise} a \pounds 2 million capital expenditure. Legitimation (the organisational mechanism by which decisions acquire binding authority) is a function of hierarchy that AI does not replace. A judgment broker who validates a shadow workflow is exercising legitimate authority derived from their position in the organisational structure. Remove the structure, and the authority evaporates.

\subsection{Accountability}

Who is fired when things go wrong? ``The mesh decided'' is not an answer a board of directors accepts. Accountability requires identifiable humans who bear consequences for outcomes. The judgment broker role addresses this partially (the broker who validates a workflow bears accountability for that validation), but the full accountability architecture for agentic systems is not worked out. When an agent in the mesh makes a decision that produces a regulatory violation, the chain of accountability must be traceable to a human. The declarative governance layer (Chapter~\ref{ch:governance}) provides audit trails, but audit trails are not the same as accountability structures.

\subsection{Political Negotiation}

Whose priorities win when budgets are finite? AI can optimise within constraints; it cannot set the constraints. The allocation of scarce resources across competing priorities is fundamentally a political act: a negotiation among stakeholders with different interests, values, and power. The mesh can surface information that makes the negotiation more transparent, but it cannot replace the negotiation itself.

\begin{cautionbox}[title={An Honest Reframing}]
If three of hierarchy's five core functions persist (legitimation, accountability, and political negotiation), then ``coordination collapse'' may be an overstatement. A more honest framing: \textbf{information routing collapse} --- one function of hierarchy dissolves, forcing the others to restructure around different assumptions. The coordination tax shrinks dramatically; it does not reach zero. This report uses the more provocative framing because it captures the magnitude of the disruption, but the qualification belongs in the analysis.
\end{cautionbox}

\textbf{Research needed}: How do legitimation, accountability, and political negotiation operate in post-hierarchical structures? Haier's micro-enterprise model provides the best existing evidence but operates in a culturally specific context that may not generalise.

\section{Evidence Base Weaknesses}

The argument in this report rests on specific empirical findings. Several of these findings have limitations that deserve explicit acknowledgement.

\subsection{The Jagged Frontier Map Is Outdated}

The BCG experiment \parencite{DellAcqua2026} used GPT-4 (now two model generations behind current frontier models). The specific task boundaries identified in that study --- the 18 tasks inside the frontier and the business judgment task outside it --- have almost certainly shifted. Models released since the study likely perform well on tasks that were outside GPT-4's frontier.

This does not invalidate the concept. The jagged frontier is a structural feature of AI capability, not a fixed map. The principle (that adjacent tasks can fall on opposite sides of the boundary, and that the boundary is unpredictable) holds regardless of where the boundary currently sits. But any practitioner using this report to make operational decisions must retest the frontier with current models in their specific domain. The 2023 map is a historical reference, not a current navigation chart.

The concern is now partially answered by continuous measurement rather than a fixed retest. The Remote Labor Index \parencite{RLI2025, CAIS2026RLI} provides exactly the longitudinal, client-quality re-measurement this section called for: automation of real freelance deliverables (judged against a paying-client bar, not a benchmark proxy) rose from 2.5 to 16.1 per cent in eight months. That trajectory confirms both halves of the jagged frontier claim simultaneously: the frontier is still jagged (professional-quality performance still fails on the large majority of projects) and it is moving fast (the boundary shifted 6.4-fold in under a year). The residual gap is real, though: the RLI measures freelance deliverables, not the coordination tasks inside an organisation that this report is actually about. A practitioner still has to map their own domain's frontier. The RLI tells them how urgently, not where.

\subsection{The 17$\times$ Error Source: Resolved}

The first edition of this report leaned on a 17-times error multiplication statistic \parencite{TDS2026} to ground the case against naive multi-agent deployment, support the judgment broker role, and supply the quantitative punch behind the ant death spiral narrative. The source was Towards Data Science --- a community-edited blogging platform, not a peer-reviewed journal. This chapter flagged it as the report's weakest load-bearing number. The literature has since delivered a replacement.

\textcite{Cemri2025}, published at NeurIPS 2025 (Datasets and Benchmarks Track) and therefore peer-reviewed, documents failure rates of 41 to 86.7 per cent across seven state-of-the-art open-source multi-agent frameworks (ChatDev, MetaGPT, HyperAgent, AppWorld, AG2, Magentic-One, and OpenManus), drawn from more than 1,600 annotated execution traces with inter-annotator agreement of Cohen's $\kappa = 0.88$. The paper does more than replace a number: it supplies a causal taxonomy of fourteen failure modes across three categories (specification issues, inter-agent misalignment, and task verification), so the error-multiplication argument now rests on a documented mechanism rather than an unexplained multiplier. \parencite{Khanal2026} adds an independent second line of evidence: across 23,392 episodes, frontier models show the \emph{highest} long-horizon meltdown rates, up to 19 per cent, precisely because their capacity for ambitious multi-step strategies is what lets them spiral further before failing. Memory scaffolds, the intuitive fix, universally hurt long-horizon performance across every model tested.

\begin{cautionbox}[title={What Changed}]
The phenomenon the 17$\times$ figure gestured at (error propagation and amplification in multi-agent chains) was always real. What was weak was the citation, not the claim. \parencite{Cemri2025} and \parencite{Khanal2026} now give that claim a peer-reviewed evidence base, a causal taxonomy, and a falsification of the ``more memory fixes it'' intuition. The 17$\times$ figure is retired from this edition and survives only as the superseded early estimate it always was.
\end{cautionbox}

\subsection{The Klarna Narrative Is Cherry-Picked}

Klarna's 2025 financials showed improved margins alongside the customer satisfaction decline described in Chapter~\ref{ch:three-models}. The CEO's admission that ``cost unfortunately seems to have been a too predominant evaluation factor'' \parencite{Klarna2025} comes from a longer statement that also defended the AI strategy as a necessary evolution. Presenting only the quality degradation without the financial improvement is selective evidence: the kind of rhetorical move that a hostile reader will identify and use to dismiss the broader argument.

The honest presentation: Klarna achieved short-term cost reduction, experienced quality degradation, began rehiring humans, and the CEO acknowledged the trade-off. The lesson is not ``AI customer service fails'' but ``optimising for cost alone, without maintaining human oversight at quality-sensitive boundaries, produces predictable quality collapse.'' This is a more nuanced and more defensible claim.

The July 2026 picture sharpens this lesson rather than complicating it. Klarna's current operating model is not a simple reversal to the pre-2023 staffing level but an explicit AI-plus-human hybrid \parencite{KlarnaHybrid2026}: AI continues to handle high-volume routine queries, while human staff (hired on flexible, gig-style contracts rather than reinstated as the original full-time roles) are reserved for the empathy, discretion, and escalation cases the CEO's own admission identified as under-served. This is not evidence against AI customer service; it is evidence against one specific configuration of it. The stable end state is neither ``all AI'' nor ``all human'' but a designed division of labour at the jagged frontier --- the judgment-broker architecture this report proposes, arrived at independently by a company that started from the opposite assumption.

\subsection{The 89 Per Cent Reduction}

The claim that providing approved AI tools reduces unauthorised use by 89 per cent is used repeatedly in this report. No primary source is cited. The statistic is plausible (it aligns with shadow IT research showing that sanctioned alternatives reduce unsanctioned adoption), but a claim used this frequently requires a traceable citation. This is an open research gap.

\subsection{Vendor Report Methodology}

The shadow AI statistics used throughout the report stack vendor-funded research (Gartner, IBM, Capgemini) that has methodological limitations and commercial incentives to inflate threat numbers. No error bars, confidence intervals, or methodological caveats are provided. A reader from an academic background will note this immediately. The statistics should be treated as directionally correct rather than precisely accurate.

\section{Audience Blind Spots}

\subsection{No Day-in-the-Life}

Middle managers will ask: ``What does my Monday morning look like as a judgment broker? What am I literally doing at 9am?'' This report offers the conductor metaphor, the role transformation table, the three irreplaceable capacities --- yet it provides no concrete job specification, no competency framework, no salary benchmarking, no day-in-the-life narrative. The abstraction is intellectually satisfying but operationally insufficient for a middle manager trying to decide whether to invest in this transition.

\textbf{Research direction}: develop detailed judgment broker role specifications for three to five industry verticals, including day-in-the-life scenarios, competency models, and compensation benchmarking.

\subsection{No Total Cost of Ownership}

AI leaders will ask: ``What does this cost for a 500-person organisation? What is the ROI timeline?'' Chapter~\ref{ch:implementation} provides an indicative cost tier model, but this falls short of a total cost of ownership analysis that includes ongoing compute, model API costs (which are volatile and trending downward), knowledge engineering maintenance, judgment broker training, and the opportunity cost of the transition period. Without financial modelling, the implementation roadmap is a strategic argument, not a business case.

\textbf{Research direction}: develop a TCO calculator parameterised by organisation size, industry vertical, regulatory burden, and existing technology infrastructure.

\subsection{The Roadmap Contradiction}

Change managers will notice that this report attacks Kotter and ADKAR for being phased transformation models, then offers a phased 90-day roadmap. Chapter~\ref{ch:change-problem} addresses this contradiction explicitly (framing the roadmap as a bootstrap sequence rather than a change plan), but the resolution requires the reader to accept a distinction that may feel like sophistry. The tension is genuine and deserves continued intellectual engagement rather than a single paragraph of resolution.

\subsection{The Analyst-Vendor Conflict}

All three audiences will notice that DreamLab AI Consulting is simultaneously the analyst (producing this report), the framework creator (proposing the Dynamic Agentic Mesh), and the vendor (operating VisionClaw). This is a classic consulting positioning conflict. The report's credibility depends on acknowledging this conflict directly: ``We built this because we believe it. Here is the evidence for and against. Judge the framework on its merits, not its provenance.''

\section{Engaging with Acemoglu}

Daron Acemoglu's work on AI and productivity has been cited in this report's bibliography but not substantively engaged with in the main text. This is a significant omission because Acemoglu's thesis directly challenges the optimistic productivity assumptions underlying the coordination collapse argument.

\begin{evidencebox}[title={The Acemoglu Sceptical Thesis}]
In ``The Simple Macroeconomics of AI'' \parencite{Acemoglu2025}, Acemoglu argues that AI's impact on total factor productivity (TFP) will be modest (at most \textbf{0.66 per cent over 10 years}. His reasoning:

\begin{enumerate}
  \item AI automates tasks, not jobs. The tasks most amenable to automation are not necessarily the tasks that contribute most to productivity.
  \item The ``easy-to-learn, hard-to-verify'' problem: many tasks that AI can perform require costly human verification, limiting net productivity gains. This maps directly onto the vigilance problem described in Chapter~\ref{ch:vigilance}.
  \item New task creation (historically the mechanism by which automation generates employment growth) may be slower with AI than with previous technologies.
  \item AI may increase inequality rather than aggregate productivity, concentrating gains among capital owners and high-skill workers while displacing middle-skill employment.
\end{enumerate}
\end{evidencebox}

The coordination collapse thesis should engage with Acemoglu on three fronts.

\textbf{First}, the 0.66 per cent TFP estimate applies to the economy as a whole, not to individual organisations. An organisation that deploys an agentic mesh effectively may capture disproportionate productivity gains precisely because its competitors do not. The compounding organisation thesis does not require AI to transform the macroeconomy; it requires AI to provide a competitive advantage to organisations that use it well. This is consistent with Acemoglu's framework: the gains are real but unevenly distributed.

\textbf{Second}, the task-level analysis is compatible with the jagged frontier concept. Acemoglu's point that AI automates tasks, not jobs, aligns with the argument that the judgment broker's role is to determine which tasks sit inside the frontier and which require human judgment. The mesh architecture is explicitly designed around task decomposition, not wholesale job replacement. Where Acemoglu warns that the ``easy-to-learn, hard-to-verify'' problem limits productivity gains, the mesh responds with HITL Precision and Trust Variance monitoring (architectural mechanisms for managing verification costs).

\textbf{Third}, the equity concern is serious and currently unaddressed in this report. If the mesh concentrates productivity gains among organisations with the capital and technical expertise to deploy it, while displacing middle-skill workers in organisations that cannot, the net social effect may be increased inequality even as individual organisations benefit. This is a genuine tension that the coordination collapse thesis does not resolve.

\textbf{Fourth}, the RLI trajectory \parencite{RLI2025, CAIS2026RLI} puts pressure on Acemoglu's implicit pace assumption without threatening his mechanism. An automation rate that moved from 2.5 to 16.1 per cent of client-quality freelance deliverables in eight months is a fast-moving frontier by any standard; if that pace continued, the case for a durable 0.66 per cent TFP ceiling would look increasingly strained. But 16.1 per cent automation of freelance deliverables is entirely compatible with modest aggregate productivity gains if verification and integration costs dominate. Acemoglu's ``easy-to-learn, hard-to-verify'' problem, restated at the level of task completion rather than macroeconomic output, maps directly onto this constraint. \textcite{Liu2026} now gives that mechanism a formal home: contextual transaction cost includes an explicit verification-burden term, and the finding that collective efficiency collapses 91.4 per cent from the lowest to the highest CTC quartile is, in effect, a microfoundation for why task-level automation does not translate one-for-one into organisational or macroeconomic productivity (Chapter~\ref{ch:coordination-economics}).

\textbf{Fifth}, the firm-level evidence assembled since April supports both halves of this report's reading, and Acemoglu's distributional concern besides. \textcite{RampRevelio2026} and the Box survey reported by \textcite{Levie2026} (the latter resting on a single social-media source without published methodology) show real gains at the firm level, but gains concentrated among organisations that have already adopted AI intensively, with high-intensity adopters growing headcount 10.2 per cent over two years while low-intensity adopters showed no significant change. That is uneven distribution \emph{across} organisations, exactly as this report's compounding-organisation argument predicts. \parencite{Brynjolfsson2025Canaries} supplies the complementary, and more troubling, half: uneven distribution \emph{within} the labour market, with early-career workers in AI-exposed occupations showing a roughly 16 per cent relative employment decline since late 2022. Acemoglu's equity concern was speculative in the first edition of this report. It is now measured.

\section{New Open Questions (July 2026)}

The evidence that arrived since April resolves several of this chapter's original flags. It also opens new ones. Five deserve explicit tracking.

\begin{itemize}
  \item \textbf{CTC measurement in production.} \textcite{Liu2026} gives contextual transaction cost a rigorous definition and shows it predicts collective efficiency in simulation and in trace-instrumented real-LLM runs. The instrumentation is no longer absent: AgentBox now runs a contextual-transaction-cost emitter on its hook spine, forwarding a per-workflow token-count-and-handoff envelope, above a native context-compression layer (\texttt{headroom-napi}/CCR) whose audit trail is barred by hard-coded rule from ever being compressed (Chapter~\ref{ch:agentbox}). What genuinely remains open is comparability, not existence: no agreed unit or method yet lets one mesh's CTC be read against another's, and the first production figure from that emitter is still to be reported (Recommendation~3 of Chapter~\ref{ch:ecosystem-roadmap}). The KPIs proposed in Chapter~\ref{ch:new-kpis} now have both a theoretical anchor and a running meter; what they lack is a cross-mesh baseline against which a given reading counts as good or bad.
  \item \textbf{The broker-apprenticeship problem.} If entry-level roles in AI-exposed occupations decline \parencite{Brynjolfsson2025Canaries} while adopting firms hire AI-fluent juniors rather than training them from scratch \parencite{RampRevelio2026}, where does the next generation of judgment brokers learn the domain expertise the role requires? The answer plausibly differs for adopting and non-adopting firms, which would make broker scarcity a compounding advantage for early movers. This is a mechanism this report has not modelled.
  \item \textbf{The design question is contested.} \textcite{Dochkina2026} reports self-organising agent collectives, with autonomous role selection, outperforming both fully centralised coordination and fully designed structures. This contrasts with the finding of \textcite{Liu2026} that adaptive, CTC-aware orchestration beats unstructured collectives. Both are single studies with different task distributions and scales. The field has not converged on whether agent collectives should be designed, allowed to self-organise, or some specific hybrid of the two, and this report should not pretend otherwise.
  \item \textbf{Survivor bias in the deployment evidence.} \textcite{Pereira2026Playbook} study the successful 5 per cent of deployments against the baseline of \textcite{MITNANDA2025}, in which most pilots produced no measurable P\&L impact. That comparison is valuable for identifying what successful deployments have in common, but it cannot answer how common success is, because the sample is selected on the outcome. Prevalence estimates should not be derived from playbook-style evidence, however good the playbook.
  \item \textbf{The RLI-to-jobs mapping.} The Remote Labor Index measures automation at the level of individual, client-quality deliverables. No method yet exists to translate a deliverable-level automation rate into an occupation-level displacement forecast. The tasks-versus-jobs distinction that \textcite{Chatterji2026} raises applies with full force to the RLI's own headline number. Sixteen per cent of freelance deliverables being automatable at professional quality is not the same claim as sixteen per cent of any job disappearing, and this report has not built the bridge between the two.
\end{itemize}

\section{Energy and Environmental Costs}

The report proposes deploying agentic infrastructure at organisational scale without addressing the energy implications. This is an increasingly visible omission as the environmental costs of AI infrastructure become a matter of public and regulatory concern.

\begin{evidencebox}[title={The Energy Cost of Agentic Systems}]
\begin{itemize}
  \item \textcite{Patterson2025} at Microsoft Research find that reasoning-intensive queries (the kind that agentic orchestration relies on) have a median energy cost of 4.32 Wh per query, compared to 0.34 Wh for standard queries. This represents a \textbf{43$\times$ increase} in energy consumption for the reasoning-heavy workloads that an agentic mesh generates continuously.
  \item \textcite{NatureSustainability2025} projects that US AI server infrastructure will produce \textbf{24--44 Mt CO$_2$-equivalent annually} through 2030. This range is growing as model sizes increase and inference demand scales.
  \item An agentic mesh running continuous orchestration, discovery monitoring, and governance checks generates substantially more inference load than on-demand AI tool use. The compound flywheel, by design, \emph{increases} inference volume as the mesh becomes more active.
\end{itemize}
\end{evidencebox}

The energy cost of the mesh is not merely a corporate social responsibility issue. It is a business cost that scales with mesh activity. Organisations deploying agentic infrastructure must account for compute costs that are qualitatively different from traditional IT expenditure: costs that increase as the mesh becomes more capable, not less.

\textbf{Research direction}: develop energy cost models for agentic mesh deployment at different organisational scales, including comparison with the energy cost of the human coordination infrastructure the mesh replaces (office buildings, commuting, business travel). The net energy equation is not obviously negative --- a 200-person management layer has its own carbon footprint --- and it is not obviously positive.

\section{Cybersecurity Attack Surface}

The statistic that financial services organisations manage 96 machine identities per human \parencite{CyberArk2025} (with only 31 per cent having controls in place) is cited in Chapter~\ref{ch:industry-landscape} as a governance issue. It is also a security issue that the report does not adequately address. An agentic mesh creates a dramatically expanded attack surface:

\begin{itemize}
  \item Each agent skill is a potential vector for prompt injection: an adversary who can manipulate an agent's input can redirect its behaviour within its authority boundary.
  \item The OWL~2 ontology that provides shared semantics is a single point of corruption: poisoning the ontology could redirect agent routing across the entire mesh.
  \item The Nostr identity layer provides cryptographic verification but depends on key management practices that most organisations have not developed for non-human identities.
  \item The Insight Ingestion Loop creates a pathway from frontline input to production workflows: a supply chain vulnerability that adversaries could exploit by planting malicious ``shadow workflows.''
\end{itemize}

\textbf{Research direction}: systematic threat modelling for agentic mesh architectures, including adversarial analysis of the Insight Ingestion Loop and ontology poisoning scenarios.

\section{Equity and Access}

Who benefits from the mesh and who is structurally excluded? The report addresses middle managers (a specific professional cohort) but does not address:

\begin{itemize}
  \item Workers without AI literacy who cannot participate in the discovery layer.
  \item Roles that do not lend themselves to agent augmentation: physical labour, care work, creative work that resists decomposition into tasks.
  \item Organisations in the Global South that lack the infrastructure and capital to deploy agentic systems.
  \item The distributional effects within organisations: does the mesh flatten compensation as it flattens hierarchy, or does it concentrate value in the judgment broker tier?
\end{itemize}

Acemoglu's concern about AI increasing inequality rather than productivity is most acute here. If the Dynamic Agentic Mesh is a tool available primarily to well-capitalised organisations in knowledge-intensive industries, its net effect may be to widen the gap between those organisations and everyone else. This is a dynamic that the \textcite{WEF2025} Future of Jobs Report anticipates when it projects 170 million new roles created alongside 92 million displaced.

\section{Competing Frameworks Not Addressed}

The report positions against Block and Every. Neither is a consulting framework the target audience is actively evaluating. It does not differentiate substantively from:

\begin{itemize}
  \item \textbf{SAFe 6.0} \parencite{SAFe2024}, which added AI-specific agile change patterns in 2024.
  \item \textbf{Microsoft Copilot Studio} \parencite{MicrosoftCopilotStudio2025}, an enterprise organisational AI deployment framework with the largest existing footprint.
  \item \textbf{Holacracy 5.0} \parencite{Holacracy2024}, which specifically addressed the scale failures attributed to earlier versions.
  \item \textbf{Anthropic's constitutional AI governance}, directly relevant to the governance paradox in Chapter~\ref{ch:governance}.
\end{itemize}

A competitor comparison table is provided in Appendix~\ref{ch:competitor-matrix}, but the main text does not engage with these alternatives at the depth they deserve.

\section{The VisionClaw Credibility Gap}

Chapter~\ref{ch:technical-substrate} describes VisionClaw honestly, including both its genuine capabilities and its limitations. But the credibility gap between a 15-person creative technology team and a 500-person manufacturing firm has not been bridged. The favourable conditions under which VisionClaw operates (high technical literacy, low process rigidity, flat power dynamics, intrinsic motivation, the system architect as the primary user) are the most favourable possible environment.

\begin{keyinsight}
VisionClaw is an existence proof under ideal conditions, not a proof of general applicability. The coordination collapse thesis requires bridging evidence that the architecture works in less favourable environments: regulated industries, non-technical workforces, unionised labour, legacy ERP systems, and middle managers who did not choose the system.

Four bridging strategies are available:
\begin{enumerate}
  \item Run a pilot in a genuinely different context (regulated industry, non-technical workforce).
  \item Partner with an established consulting framework (Deloitte, BCG) for independent validation.
  \item Publish the methodology as open-source and invite external testing.
  \item Position VisionClaw explicitly as ``existence proof under favourable conditions'' and build the evidence base incrementally.
\end{enumerate}
The fourth strategy is the one this report adopts. The others represent the research agenda.
\end{keyinsight}

\section{Additional Failure Modes}

The report documents two failure modes in detail: the ant death spiral (agents responding to agents in infinite loops) and the vigilance degradation problem (humans falling asleep at the wheel as AI quality improves). But an agentic mesh at organisational scale introduces additional failure modes that this analysis does not adequately address:

\begin{itemize}
  \item \textbf{Cascading errors}: a single incorrect agent output propagating through downstream agents before the judgment broker can intervene.
  \item \textbf{Single point of failure in the judgment broker}: if the broker is unavailable, incapacitated, or overwhelmed, the mesh has no fallback governance mechanism.
  \item \textbf{Ontology drift}: the OWL~2 ontology gradually diverging from the actual business domain as the organisation evolves, producing increasingly inaccurate agent routing.
  \item \textbf{Metric gaming}: agents optimising for the four KPIs in ways that satisfy the metrics while degrading actual performance. This is Goodhart's law applied to agentic governance.
  \item \textbf{Cultural rejection}: the mesh may function technically while being socially rejected by the workforce, producing compliance without engagement.
\end{itemize}

\textbf{Research direction}: systematic failure mode analysis (FMEA) for agentic mesh architectures, with probability and severity estimates calibrated to organisational context.
